% GENERATED FILE. Edit canonical lessons, then run npm run generate:books.
% canonical-source-hash: fnv1a64:90513703f7ec0d30
% canonical-lessons: GE-C261-ich-haette-gern, GE-C261-ich-wuerde, GE-C261-wenn-ich-zeit-haette, GE-C261-review, GE-W261-sechs-zeilen

\chapter{Would — würde, hätte, wäre, and Six Lines}
\label{ch:ge-would}
\hlchaptermodality{261}

\begin{chapteropening}
\textbf{By the end of this chapter:} \emph{I can say what would happen with würde, hätte, wäre and könnte, and write six connected lines about a trip that end on a condition.}

Complete GE-W261-sechs-zeilen: six lines to one reader about a trip, events in order in the past, a reason, a feeling, and a last line with wenn and the Konjunktiv II.
\end{chapteropening}

\section[ich hätte gern / das wäre]{ich hätte gern — I'd like, and that would be}

\label{lesson:GE-C261-ich-haette-gern}



\begin{quote}
Say \textquotedblleft{}I had\textquotedblright{} and \textquotedblleft{}I was\textquotedblright{}. (\emph{Ich hatte. Ich war.}) Change the \emph{a} of each to \emph{ä}, add an \emph{e} to the second, and the past turns into \textquotedblleft{}would\textquotedblright{}.
\end{quote}

\subsection*{You'll want to know: ich hätte gern}
\textbf{ich hätte gern} — \textquotedblleft{}I'd like\textquotedblright{}, literally \textquotedblleft{}I would gladly have\textquotedblright{}: \emph{ikh HET-tuh gairn}.

\begin{quote}
\emph{Ich hätte gern einen Kaffee.} — I'd like a coffee. \emph{Ich hätte gern eine Fahrkarte nach Berlin.} — I'd like a ticket to Berlin.
\end{quote}

It sits beside \emph{ich möchte}, and both are polite. At a counter or in a café, \emph{ich hätte gern} is a common way to order.

The forms: \emph{ich hätte}, \emph{du hättest}, \emph{er hätte}, \emph{wir hätten}, \emph{Sie hätten}.

\begin{grammarlens}[title={Grammar Lens: das wäre}]
\textbf{wäre} is \textquotedblleft{}would be\textquotedblright{}, made from \emph{war}:

\begin{quote}
\emph{Das wäre schön.} — That would be lovely. \emph{Das wäre nett.} — That would be kind.
\end{quote}

The forms: \emph{ich wäre}, \emph{du wärst}, \emph{er wäre}, \emph{wir wären}, \emph{Sie wären}.

This is the \textbf{Konjunktiv II}, the mood of what is not real, or not yet, or said gently. You have used one form already: \emph{du solltest}, which is spelled like the Präteritum. \emph{Hätte} and \emph{wäre} add an umlaut to theirs.
\end{grammarlens}

\begin{cousinweb}
\textbf{hätte} is \emph{hatte} and \textbf{wäre} is \emph{war}, each with an umlaut. Neither form is in EtymDB 2.1, so this is the usual account in histories of German: the old subjunctive endings had an \emph{i} in them, as in Old High German \emph{wārī}, and that \emph{i} pulled the vowel forward, \emph{a} to \emph{ä}. The \emph{i} is gone; the umlaut it left is the mark of the mood. \emph{Haben} itself EtymDB traces to Proto-Germanic \emph{*habjaną}.
\end{cousinweb}

\subsection*{Your turn}
Say these aloud:

\begin{itemize}
\raggedright
  \item \emph{hatte}, \emph{hätte}; \emph{war}, \emph{wäre}
  \item \emph{Ich hätte gern einen Kaffee.}
  \item \emph{Ich hätte gern eine Fahrkarte nach Berlin, hin und zurück.}
  \item \emph{Das wäre schön.}
  \item \emph{Das wäre nett.}
  \item a problem, then a polite request: \emph{Ich habe meine Fahrkarte verloren. Ich hätte gern eine neue.}
\end{itemize}

\begin{tcolorbox}[breakable,colback=teal!4,colframe=teal!35!black,title={Before you move on}]
\textquotedblleft{}I'd like a coffee\textquotedblright{}? (\emph{\textbf{Ich hätte gern einen Kaffee.}}) \textquotedblleft{}That would be lovely\textquotedblright{}? (\emph{\textbf{Das wäre schön.}}) What turns \emph{hatte} into \emph{hätte}? (\textbf{An umlaut.})
\end{tcolorbox}
\section[ich würde + infinitive]{ich würde — would}

\label{lesson:GE-C261-ich-wuerde}



\begin{quote}
Say \textquotedblleft{}I'd like a coffee\textquotedblright{} and \textquotedblleft{}that would be lovely\textquotedblright{}. (\emph{Ich hätte gern einen Kaffee. Das wäre schön.})

\emph{Haben} and \emph{sein} have \textquotedblleft{}would\textquotedblright{} forms of their own. Most verbs borrow one.
\end{quote}

\begin{grammarlens}[title={Grammar Lens: würde, and the infinitive last}]
For \textquotedblleft{}would\textquotedblright{} with an ordinary verb, German takes \textbf{würde} and puts the verb's infinitive at the end, the same bracket a modal builds:

\begin{quote}
\emph{Ich würde gern nach Italien} \textbf{reisen}. — I would like to travel to Italy. \emph{Was würdest du} \textbf{machen}? — What would you do?
\end{quote}

The forms: \emph{ich würde}, \emph{du würdest}, \emph{er würde}, \emph{wir würden}, \emph{ihr würdet}, \emph{Sie würden}.

After \emph{dass} or \emph{weil}, \emph{würde} goes to the very end, after the infinitive:

\begin{quote}
\emph{Ich glaube, dass er gern} \textbf{kommen würde}.
\end{quote}

In everyday German the split is roughly this: \emph{haben}, \emph{sein} and the modals keep their one-word forms, \emph{hätte}, \emph{wäre}, \emph{sollte}; most other verbs take \emph{würde}.
\end{grammarlens}

\begin{cousinweb}
\textbf{würde} is \textbf{wurde}, the Präteritum of \emph{werden}, \textquotedblleft{}to become\textquotedblright{}, with an umlaut. That is the same move as \emph{hatte} to \emph{hätte} and \emph{war} to \emph{wäre}: an umlaut on the past marks the Konjunktiv II. Not every verb takes one: \emph{sollte}, \emph{wollte} and the regular \emph{-te} verbs look exactly like their past, and the sentence tells you which is meant.

\emph{Werden} $\leftarrow$ Old High German \emph{werdan} $\leftarrow$ Proto-Germanic \emph{*werþaną}, \textquotedblleft{}to turn\textquotedblright{} (EtymDB 2.1). EtymDB does not list \emph{würde}; the umlaut history is the usual account.
\end{cousinweb}

\subsection*{Your turn}
Say these aloud:

\begin{itemize}
\raggedright
  \item \emph{wurde}, then \emph{würde}
  \item \emph{Ich würde gern nach Italien reisen.}
  \item \emph{Was würdest du machen?}
  \item \emph{Ich glaube, dass er gern kommen würde.}
  \item three pasts, three umlauts: \emph{hatte, hätte; war, wäre; wurde, würde}
  \item an honest answer: \emph{Was würdest du machen? Das hängt davon ab.}
\end{itemize}

\begin{tcolorbox}[breakable,colback=teal!4,colframe=teal!35!black,title={Before you move on}]
\textquotedblleft{}I would like to travel\textquotedblright{}? (\emph{\textbf{Ich würde gern reisen.}}) Where does the infinitive go? (\textbf{To the end.}) Which past is \emph{würde} made from? (\emph{\textbf{Wurde}}, from \emph{werden}.)
\end{tcolorbox}
\section[wenn ich Zeit hätte, würde ich …]{wenn ich Zeit hätte — if things were different}

\label{lesson:GE-C261-wenn-ich-zeit-haette}



\begin{quote}
Say \textquotedblleft{}if I have time, I'll come\textquotedblright{}. (\emph{Wenn ich Zeit habe, komme ich.}) That condition may well come true. Now one that is not true at all.
\end{quote}

\begin{grammarlens}[title={Grammar Lens: the Konjunktiv II in both halves}]
When the condition is not real, both halves go into the Konjunktiv II:

\begin{quote}
\emph{Wenn ich Zeit} \textbf{hätte}, \textbf{würde} \emph{ich mehr lesen.} — If I had time, I would read more.
\end{quote}

I do not have time; that is the point. Set the two side by side:

\begin{quote}
\emph{Wenn ich Zeit habe, lese ich.} — real: it may happen. \emph{Wenn ich Zeit hätte, würde ich lesen.} — not real: it is not happening.
\end{quote}

The word order has not changed: the verb at the end of the \emph{wenn} clause, then verb, comma, verb: \emph{hätte, würde}.

One short version is worth having on its own: \emph{Wenn ich du wäre …}, if I were you.
\end{grammarlens}

\begin{culture}
\textbf{ich könnte} is \textquotedblleft{}I could\textquotedblright{}, in the sense of \textquotedblleft{}I would be able to\textquotedblright{}. It is \emph{konnte} with an umlaut, and like \emph{würde} it puts the infinitive last:

\begin{quote}
\emph{Wenn ich Geld hätte, könnte ich nach Japan reisen.} — If I had money, I could travel to Japan.
\end{quote}

It also makes a suggestion soft: \emph{Wir könnten morgen fahren}, we could go tomorrow.

Listen for the vowel. \emph{Ich konnte} is \textquotedblleft{}I was able to\textquotedblright{}, over and done. \emph{Ich könnte} is \textquotedblleft{}I could\textquotedblright{}, if things were different.
\end{culture}

\begin{cousinweb}
\textbf{könnte} is \textbf{konnte} with the umlaut of the Konjunktiv II. EtymDB 2.1 traces \emph{können} through Old High German \emph{kunnan} to Proto-Germanic \emph{*kunnaną}; you met its link to English \emph{could}, and its old sense \textquotedblleft{}to know how\textquotedblright{}, with \emph{konnte} in chapter twenty-five. \textbf{Zeit} goes back to Proto-Germanic \emph{*tīdiz}, \textquotedblleft{}time, period\textquotedblright{} (EtymDB 2.1), the same word as English \emph{tide}.
\end{cousinweb}

\subsection*{Your turn}
Say these aloud:

\begin{itemize}
\raggedright
  \item \emph{Wenn ich Zeit hätte, würde ich mehr lesen.}
  \item real, then not real: \emph{Wenn ich Zeit habe, lese ich. Wenn ich Zeit hätte, würde ich lesen.}
  \item \emph{Wenn ich Geld hätte, könnte ich nach Japan reisen.}
  \item \emph{ich konnte}, then \emph{ich könnte}
  \item \emph{Wenn ich du wäre, würde ich ein Taxi nehmen.}
\end{itemize}

\begin{tcolorbox}[breakable,colback=teal!4,colframe=teal!35!black,title={Before you move on}]
\textquotedblleft{}If I had time, I would read more\textquotedblright{}? (\emph{\textbf{Wenn ich Zeit hätte, würde ich mehr lesen.}}) \textquotedblleft{}I could\textquotedblright{}, if things were different? (\emph{\textbf{Ich könnte.}}) Which is over and done, \emph{konnte} or \emph{könnte}? (\emph{\textbf{Konnte.}})
\end{tcolorbox}
\section[(review)]{Review — what would happen}

\label{lesson:GE-C261-review}



\begin{quote}
Order a coffee, politely. (\emph{Ich hätte gern einen Kaffee.}) And say \textquotedblleft{}that would be lovely\textquotedblright{}. (\emph{Das wäre schön.})
\end{quote}

\begin{grammarlens}[title={What you've built}]
Three lessons, and German can now say what is not so. Four pasts with an umlaut: \emph{hätte}, \emph{wäre}, \emph{würde}, \emph{könnte}. \emph{Würde} with the infinitive last for most verbs. And two kinds of \emph{wenn}: the present for what may well happen, the Konjunktiv II for what is not true now.

The two side by side, about the same weekend:

\begin{quote}
\emph{Wenn ich Zeit habe, fahre ich nach Berlin.} \emph{Wenn ich Zeit hätte, würde ich nach Berlin fahren.}
\end{quote}

The first is a plan. The second is a wish.
\end{grammarlens}

\subsection*{Your turn}
Say these aloud:

\begin{itemize}
\raggedright
  \item \emph{Wenn ich Zeit habe, fahre ich nach Berlin.}
  \item \emph{Wenn ich Zeit hätte, würde ich nach Berlin fahren.}
  \item \emph{Wenn ich Geld hätte, könnte ich nach Japan reisen.}
  \item \emph{Das wäre schön, aber ich habe keine Zeit.}
  \item one wish of your own: \emph{Wenn ich …, würde ich …}
  \item \emph{Das Taxi wäre schneller als die U-Bahn, aber nicht so billig wie die U-Bahn.}
  \item \emph{Vor einem Jahr war ich in Berlin. Ich lerne seit einem Jahr Deutsch. Wenn ich mehr Zeit hätte, würde ich wieder nach Berlin fahren.}
  \item \emph{Kommst du mit? Es kommt darauf an. Das hängt davon ab, ob ich Zeit habe.}
\end{itemize}

\begin{tcolorbox}[breakable,colback=teal!4,colframe=teal!35!black,title={Before you move on}]
A plan or a wish: \emph{Wenn ich Zeit hätte}? (\textbf{A wish}: it is not true now.) \textquotedblleft{}I could\textquotedblright{}? (\emph{\textbf{Ich könnte.}}) Where does the infinitive go after \emph{würde}? (\textbf{To the end.})
\end{tcolorbox}
\section[(sechs Zeilen)]{Six lines about a trip}

\label{lesson:GE-W261-sechs-zeilen}



\begin{quote}
Chapter fifty-five asked for four lines that held together, with \emph{weil}. This one asks for six, and they tell a story.

Pick one person who will read it: \emph{ein Freund}, \emph{eine Freundin}, someone in your family.
\end{quote}

\subsection*{What to know first}
A told story is more than a list of events. It has a \textbf{scene}: \emph{es war}, \emph{es regnete}. It has \textbf{events} in order, with the verb straight after each step word: \emph{zuerst}, \emph{dann}, \emph{danach}, \emph{schließlich}. It has at least one \textbf{reason}, pointing back with \emph{weil} or forward with \emph{darum}. And a good one ends by looking forward.

For a message to a friend, either past is right. The Präteritum sounds like a story being told; the Perfekt sounds like you talking: \emph{ich bin nach Berlin gefahren}. Keep \emph{war} and \emph{hatte} in the Präteritum either way.

The Goethe-Zertifikat B1 writing paper asks for personal and formal messages and a short opinion post. Six lines to a friend is a first step towards the personal message, not the whole of it.

\subsection*{Writing — connected composition}
On paper, write \textbf{six lines} to your reader about a trip, real or invented, one line for each of these:

\begin{enumerate}
\raggedright
  \item the scene: when it was, with \emph{vor} (\emph{vor einem Monat}), and the weather;
  \item the first thing that happened, with \emph{zuerst};
  \item the next thing, with \emph{dann}, and a reason, with \emph{weil} or \emph{darum};
  \item how it ended, with \emph{schließlich};
  \item how you felt, with \emph{ich war froh, dass …}, or what you thought of it, with \emph{meiner Meinung nach} and a comparison: \emph{schöner als}, \emph{so schön wie};
  \item one line forward, with \emph{wenn ich … hätte, würde ich …}, or \emph{das wäre …}, or \emph{ich hoffe, … zu …}.
\end{enumerate}

Use only words you already write. The exercise is the joining, and the place of each verb.

\begin{tcolorbox}[breakable,colback=teal!4,colframe=teal!35!black,title={Before you move on}]
Now go through your lines and find every verb. In lines two, three and four, the verb should come straight after the step word: \emph{zuerst ging ich}, not \emph{zuerst ich ging}. That is the slip to check first.

Then every \emph{weil}, \emph{dass} and \emph{wenn}. Is its verb at the end of the clause? If a line starts with the \emph{wenn} clause, does the main verb come straight after the comma: \emph{hätte, würde}?

Last, line six on its own. If it has \emph{hätte} or \emph{würde}, is it about something that is not true now? Then it looks forward, and the story is finished.
\end{tcolorbox}
